\begin{claim}{D.65}{Proved}{A uniform square obstruction and the remaining standard centralizer}
For each integer $n\geq2$, put
\[
 G_n=\langle z,v\mid v^2=1,\ [z^2,v]^n=1\rangle,\qquad
 g=z^2,\quad c=[g,v].
                                                               \tag{D.65a}
\]
These are discrete groups; $\widehat G_n$ denotes their
full profinite completions.
For an integer $r$ with $\gcd(2r+1,n)=1$, the exact
square criterion is
\[
 c^rg\text{ is a square in }\widehat G_n
       \quad\Longleftrightarrow\quad
                  r\equiv0\text{ or }-1\pmod n.              \tag{D.65b}
\]
For either allowed residue the square root already
lies in the displayed discrete group.

Consequently the full peripheral outer centralizer
of the standard one-section negative torus involution
$b_1=[\alpha_1]$, $\alpha_1(X)=Y,\alpha_1(Y)=X$, is
\[
 C_\Omega(b_1)=N_\Omega(\langle b_1\rangle)
       =\langle b_1,[\zeta_1]\rangle\cong C_2^2,\qquad
       \zeta_1(X)=X^{-1},\quad\zeta_1(Y)=Y^{-1}.
                                                               \tag{D.65c}
\]
In particular all its multiplier values are the two
global units $1,-1$. Together with D.64 this computes
both standard negative centralizers, and every fibre
of each actual standard conjugacy-orbit map has exactly
four elements, with the explicit right-coset description.
It does not assert that the two standard orbits exhaust
all negative involutions.
\end{claim}

\begin{proof}
\emph{Two exact discrete graph-of-groups presentations.}
Let
\[
 \begin{split}
 D_n&=\langle c,v\mid c^n=v^2=1,\ vcv=c^{-1}\rangle,
           \qquad w=cv,\\
 H_n&=\langle D_n,g\mid gvg^{-1}=w\rangle.
 \end{split}                                                  \tag{D.65d}
\]
The first group is the finite dihedral group of order
$2n$. Its two specified order-two subgroups are
$\langle v\rangle,\langle w\rangle$.
The second line is their HNN extension, with the
actual edge map $v\mapsto w$.
The HNN normal form embeds $D_n$ in $H_n$.
There is an actual exponent character
\[
        \ell:H_n\longrightarrow\Z,\qquad
                   \ell(g)=1,\quad\ell(c)=\ell(v)=0.          \tag{D.65e}
\]
In particular $\langle g\rangle$ is infinite cyclic.

Eliminating $c$ gives the mutually inverse presentation
maps
\[
 \begin{aligned}
 H_n&\longrightarrow
       \langle G,V\mid V^2=1,\ [G,V]^n=1\rangle,
       &(g,v,c)&\longmapsto(G,V,[G,V]),\\
 (G,V)&\longmapsto(g,v).
 \end{aligned}                                                \tag{D.65f}
\]
Indeed $gvg^{-1}=cv$ gives $[g,v]=c$.
Conversely $V[G,V]V=[G,V]^{-1}$ follows from $V^2=1$,
so the defining dihedral relation is retained.
Now form the discrete amalgam with its stated
injective edge maps:
\[
       P_n=H_n*_{\langle g\rangle=\langle z^2\rangle}
                                                 \langle z\rangle .
                                                               \tag{D.65g}
\]
Its inverse presentation maps with (D.65a) send
\[
  (g,v,c,z)\longmapsto(z^2,v,[z^2,v],z),
        \qquad (z,v)\longmapsto(z,v).
                                                               \tag{D.65h}
\]
All defining relators have just been verified in
(D.65f), and each composite fixes its generators.
Thus $P_n=G_n$ through these explicit maps.
In particular both amalgam factors embed, and $c$
has exactly order $n$.
The following global character will be used after
completion:
\[
       \vartheta:G_n\longrightarrow\Z,\qquad
                        \vartheta(z)=1,\quad\vartheta(v)=0.
                                                               \tag{D.65i}
\]
Its restriction to $H_n$ is $2\ell$.

\emph{Length-one conjugacy and its full discrete centralizer.}
In the Bass--Serre tree of (D.65d), the vertex group
is $D_n$ and the edge groups are the specified
order-two subgroups.
Both $g$ and $a_r=c^rg$ have cyclically reduced stable
letter length one. Their successive positive and
negative powers are reduced, so they translate their
axes by one edge, and both axes contain the vertex
stabilized by $D_n$.

Suppose $hgh^{-1}=a_r$ in $H_n$.
The element $h$ carries the axis of $g$ to the axis
of $a_r$. Since $a_r$ translates by one edge, some
power of $a_r$ carries the vertex $hD_n$ to $D_n$.
Multiplying $h$ by that power therefore gives a
conjugator $d\in D_n$ with the same equation.
Rearrange it as
\[
                  g d^{-1}g^{-1}=d^{-1}c^r\in D_n.
                                                               \tag{D.65j}
\]
Britton's normal form forces
$d^{-1}\in\langle v\rangle$.
There are exactly two possibilities.
For $d=1$ the equation gives $r=0\bmod n$.
For $d=v$ the exact relation is
\[
                   vgv=c^{-1}g,
                                                               \tag{D.65k}
\]
and gives $r=-1\bmod n$.
Conversely those two displayed conjugators work.
Thus
\[
           c^rg\sim_{H_n}g
          \quad\Longleftrightarrow\quad
                      r\equiv0,-1\pmod n.                   \tag{D.65l}
\]

Assume now that $r$ is neither residue.
The character (D.65e) shows that $a_r$ cannot be
$H_n$-conjugate to any element of the amalgam edge
group $\langle g\rangle$: such a conjugacy would
have to be to $g$ itself, which (D.65l) excludes.
In the tree of (D.65g), $a_r$ fixes the $H_n$ vertex
and fixes no adjacent edge. If it fixed another
vertex, the connecting path would give an adjacent
fixed edge. Hence this is its unique fixed vertex.
Every element centralizing $a_r$ must preserve that
vertex, so
\[
                        C_{G_n}(a_r)=C_{H_n}(a_r).
                                                               \tag{D.65m}
\]
This is an argument in the discrete tree; it assumes
no fixed-path assertion for a profinite tree.

For $h\in C_{H_n}(a_r)$, commuting with a nontrivial
translation preserves and orients its axis.
Again its vertices form a single orbit under the
length-one translation $a_r$.
Thus $h=a_r^m k$ for an integer $m$ and a
$k\in D_n$ centralizing $a_r$.
The same Britton reduction as in (D.65j) forces
$k\in\langle v\rangle$.
Its nonidentity possibility is excluded by the
exact dihedral calculation
\[
              a_r v a_r^{-1}=c^{\,2r+1}v\ne v,
                                                               \tag{D.65n}
\]
where the inequality uses $\gcd(2r+1,n)=1$ and
$n\geq2$.
We have proved the complete discrete equality
\[
                         C_{G_n}(a_r)=\langle a_r\rangle .
                                                               \tag{D.65o}
\]

\emph{The actual closed surface cover.}
The homomorphism
$\epsilon:G_n\to C_2$ sending both $z$ and $v$
to the nontrivial element has kernel
\[
           \Pi_n=\langle X,Y\mid [X,Y]^n=1\rangle.
                                                               \tag{D.65p}
\]
Here the exact isomorphism
$G_n\cong\Pi_n\rtimes\langle v\rangle$ has
\[
         X=zv,\quad Y=vz,\qquad
         z=Xv,\qquad vXv=Y,\quad vYv=X.
                                                               \tag{D.65q}
\]
Substitution gives $[z^2,v]=[X,Y]$.
Conversely swapping $X,Y$ preserves the relation
in (D.65p), because it inverts their commutator.
The indicated maps between the presentation and
this semidirect product therefore invert each other,
proving both the kernel and its inclusion.

Let $Q_n$ be the finite Heisenberg group
\[
 Q_n=(\Z/n)^3,\qquad
 (a,b,t)(a',b',t')=(a+a',b+b',t+t'+ab').
\]
The map
\[
  \Pi_n\longrightarrow Q_n,\qquad
    X\longmapsto(1,0,0),\quad
    Y\longmapsto(0,1,0),\quad
    [X,Y]\longmapsto(0,0,1)                                  \tag{D.65r}
\]
is onto and has a commutator of exactly order $n$.
To identify its kernel, start with the actual
degree-$n^3$ regular cover of a once-punctured torus
given by the same homomorphism from its free group.
The boundary permutation is multiplication by that
central element of order $n$.
It has exactly $n^2$ cycles, each of length $n$,
so the covering surface has $n^2$ boundary components,
each mapping with degree $n$.
Fill each by a disk.
Van Kampen kills exactly the normal closure of
the corresponding conjugates of $[X,Y]^n$.
These are exactly the Reidemeister--Schreier relators
for the kernel of (D.65r): the original normal
closure is generated by its conjugates over the
$n^3$ cosets, and the $n$ cosets in a boundary
cycle give conjugates of the same boundary relator.
Thus the resulting closed surface group is precisely
that kernel, with the inclusion supplied by the
covering construction.
Its Euler characteristic and genus are
\[
         \chi=-n^3+n^2,\qquad
                    h_n=1+\frac{n^2(n-1)}2 .
                                                               \tag{D.65s}
\]
It has index $2n^3$ in $G_n$.
This proves directly that $G_n$ is virtually a
closed surface group.

\emph{The full completion bridge.}
Martino's theorem states that every group having
a surface subgroup of finite index is conjugacy
separable \cite[Theorem 3.7 of the inspected author
version]{C-Martino2007}.
Every finite-index subgroup of $G_n$ has the same
property, by intersection with the surface group
just constructed.
Consequently $G_n$ is hereditarily conjugacy
separable, and in particular residually finite.

We give the precise centralizer consequence and
its maps. For any such group $L$ and $a\in L$,
\[
           C_{\widehat L}(a)=\overline{C_L(a)}
                         \quad\text{inside }\widehat L.       \tag{D.65t}
\]
Let $x$ belong to the left-hand side.
For a finite-index normal subgroup $N\lhd L$,
choose $h\in L$ with $x=h\widehat n$,
$\widehat n\in\overline N$.
The commuting equation gives
$h^{-1}ah=\widehat n a\widehat n^{-1}$.
Its image modulo $N$ is therefore the image of $a$,
so the actual element $h^{-1}ah$ belongs to
$L_N=N\langle a\rangle$.
Both actual elements lie in the finite-index
subgroup $L_N$, and they are
conjugate in its full completion, since
$\overline N\subseteq\widehat{L_N}$.
Conjugacy separability of $L_N$ gives a conjugator
there. Write it as $n a^k$, with $n\in N$.
It follows that
$h^{-1}ah=n a n^{-1}$, so $hn\in C_L(a)$ and
$x\in C_L(a)\overline N$.
Intersecting these neighborhoods over $N$ puts
$x$ in the right-hand side of (D.65t).
The reverse inclusion follows from the closed
commuting equation.
Finite-index subgroups carry their full induced
profinite topologies, which is the only completion
inclusion used in this proof.

Apply (D.65t) to (D.65o).
The continuous extension of (D.65i) satisfies
$\widehat\vartheta(a_r)=2$.
Hence the map
\[
       \widehat\Z\longrightarrow
             \overline{\langle a_r\rangle}
                  \subseteq\widehat G_n,\qquad t\longmapsto a_r^t
                                                               \tag{D.65u}
\]
is injective: its composite with
$\widehat\vartheta$ is multiplication by two
on $\widehat\Z$.
It is onto the indicated procyclic closure by
definition.
Thus (D.65t) is the exact full procyclic
centralizer $\langle a_r\rangle\cong\widehat\Z$,
with no missing prime coordinates.
If $q^2=a_r$ in $\widehat G_n$, then $q$
centralizes $a_r$ and equals $a_r^t$.
The equation would give $2t=1$ in $\widehat\Z$,
which is impossible modulo two.
This excludes a square for the two disallowed
residue conditions above.
For $r=0$ the element is $g=z^2$.
For $r=-1$, (D.65k) gives
$c^{-1}g=vgv=(vzv)^2$.
Since $c$ has order $n$, these formulas prove
both directions of (D.65b).

\emph{Full units and both original reflection labels.}
Return to the full standard extension
$E_1=C_2\langle v\rangle\amalg\widehat\Z\langle z\rangle$
in D.64, with $g=z^2$ and $c_0=[g,v]$.
For each $n$, the generator maps give an actual
continuous epimorphism
\[
           q_n:E_1\longrightarrow\widehat G_n,\qquad
                              z\longmapsto z,\quad v\longmapsto v.
                                                               \tag{D.65v}
\]
Their dense discrete image generates $G_n$, so
compactness makes the map onto its full completion.
Let $\beta\in C_\Omega(b_1)$ have the unique
normalized extension from D.63.
D.64 gives $C_{F_0}(v_i)=1$, so its image
$a=\Psi(g)=\Psi(z)^2$ is exactly
\[
 a=
 \begin{cases}
   c_0^{(\mu-1)/2}g,&i=0,\\
   c_0^{(\mu+1)/2}g^{-1},&i=1,
 \end{cases}
 \qquad \mu=\lambda(\beta)\in\widehat\Z^\times .
                                                               \tag{D.65w}
\]
First take $i=0$ and put $r=(\mu-1)/2\in\widehat\Z$.
Choose an integer representative $r_n$ of $r\bmod n$.
Its image in (D.65v) is $c^{r_n}g$ and is a square.
The integer $2r_n+1$ is prime to $n$, because it
is the reduction of the full unit $\mu$.
Criterion (D.65b) gives
$r_n\equiv0,-1\pmod n$.
Equivalently,
\[
                         \mu\equiv1\text{ or }-1\pmod{2n}
                             \quad\text{for every }n\geq2.
                                                               \tag{D.65x}
\]
The finite-intersection argument of D.64k now
forces a single global sign $\mu=1$ or $\mu=-1$.

For $i=1$, inversion and the specified conjugation
give the exact comparison
\[
        g^{-1}a^{-1}g
            =c_0^{-(\mu+1)/2}g
            =c_0^{((-\mu)-1)/2}g .
                                                               \tag{D.65y}
\]
This element is the square of
$g^{-1}\Psi(z)^{-1}g$.
It is therefore precisely the already treated
first-label expression for the full unit $-\mu$,
and again $\mu\in\{1,-1\}$.
No reflection label or named cusp has been rescaled.

D.64 proves that the character pair is injective
for $b_1$, and explicitly realizes all four pairs
with these two multipliers.
This gives (D.65c).
The normalizer and actual orbit-fibre conclusions
follow exactly as in D.64.
The proof restricts the standard centralizers;
the separate existence of a conjugator for an
arbitrary negative involution remains the full
equation in D.62.

The program
\path{code/part_c/standard_oneclass_square_d65.py}
checks the defining word maps, the finite
dihedral identities in the axis calculation,
and the exact Heisenberg covering data.
The square obstruction in the full completion
uses the proved discrete centralizer and the
specified hereditary-conjugacy-separability bridge,
beyond those finite checks.
\end{proof}
